\medskip

\emph{Dans cet exercice, tous les résultats seront arrondis à $10^{-3}$ près en cas de besoin.\\
Les deux parties de cet exercice sont indépendantes l'une de l'autre.}

\bigskip

\textbf{Partie A}

\medskip

Au tennis, le joueur qui est au service peut, en cas d'échec lors du premier service, servir une deuxième balle.

En match, Abel réussit son premier service dans 70\,\% des cas. Lorsque le premier service est réussi, il gagne le point dans 80\,\% des cas.

En revanche, après un échec à son premier service, Abel gagne
le point dans 45\,\% des cas.

\medskip

Abel est au service.

On considère les évènements suivants:

\begin{itemize}[label=$\bullet~~$]
\item $S$ : \og Abel réussit son premier service \fg
\item $G$ : \og Abel gagne le point \fg.
\end{itemize}

\medskip

\begin{enumerate}
\item Décrire l'évènement $S$ puis traduire la situation par un arbre pondéré.
\item Calculer $P(S \cap G)$.
\item Justifier que la probabilité de l'évènement $G$ est égale à $0,695$.
\item Abel a gagné le point. Quelle est la probabilité qu'il ait réussi son premier service ?
\item Les évènements $S$ et $G$ sont-ils indépendants ? Justifier.
\end{enumerate}

\textbf{Partie B}

\medskip

À la sortie d'une usine de fabrication de balles de tennis, une balle est jugée conforme dans 85\,\% des cas.

\medskip

\begin{enumerate}
\item On teste successivement $20$~balles. On considère que le nombre de balles est suffisamment grand pour assimiler ces tests à un tirage avec remise. On note $X$ la variable aléatoire qui compte le nombre de balles conformes parmi les $20$ testées.
	\begin{enumerate}
		\item Quelle est la loi suivie par $X$ et quels sont ses paramètres ? Justifier.
		\item Calculer $P(X \leqslant 18)$.
		\item Quelle est la probabilité qu'au moins deux balles ne soient pas conformes parmi les $20$ balles testées ?
		\item Déterminer l'espérance de $X$.
	\end{enumerate}
\item On teste maintenant $n$ balles successivement. On considère les $n$ tests comme un échantillon
de $n$ variables aléatoires $X$ indépendantes suivant la loi de Bernoulli de paramètre $0,85$.
	
On considère la variable aléatoire
	
\[M_n = \displaystyle\sum_{i=1}^n \dfrac{X_i}{n} = \dfrac{X_1}{n} + \dfrac{X_2}{n} + \dfrac{X_3}{n} + \ldots + \dfrac{X_n}{n}\]

	\begin{enumerate}
		\item Déterminer l'espérance et la variance de $M_n$.
		\item Après avoir rappelé l'inégalité de Bienaymé-Tchebychev, montrer que, pour tout entier naturel $n$,\: $P(0,75 < M_n < 0,95) \geqslant  1 - \dfrac{12,75}{n}$.
		\item En déduire un entier $n$ tel que la moyenne du nombre de balles conformes pour un échantillon de taille $n$ appartienne à l'intervalle ]0,75~;~0,95 [ avec une probabilité
supérieure à $0,9$.
	\end{enumerate}
\end{enumerate}

\vspace{0,5cm}
